\section*{Решения}
\addcontentsline{toc}{section}{Решения}

\begin{enumerate}
\item \textit{(Демо 2026, №9)}
\[
A=\begin{pmatrix}0&2\\-3&5\end{pmatrix}.
\]
Характеристический многочлен:
\[
\chi_A(\lambda)=\det(A-\lambda I)=\lambda^2-5\lambda+6=(\lambda-2)(\lambda-3).
\]
Следовательно, собственные значения: \(\lambda_1=2\), \(\lambda_2=3\).

Для \(\lambda=2\):
\[
(A-2I)x=0 \iff
\begin{pmatrix}-2&2\\-3&3\end{pmatrix}x=0
\iff x_1=x_2,
\]
поэтому
\[
E_{2}=\langle (1,1)^T\rangle.
\]

Для \(\lambda=3\):
\[
(A-3I)x=0 \iff
\begin{pmatrix}-3&2\\-3&2\end{pmatrix}x=0
\iff 3x_1=2x_2,
\]
поэтому
\[
E_{3}=\langle (2,3)^T\rangle.
\]

Матрица имеет два различных собственных значения, значит, она диагонализуема. Возьмём
\[
P=\begin{pmatrix}1&2\\1&3\end{pmatrix},\qquad
D=\begin{pmatrix}2&0\\0&3\end{pmatrix},\qquad
P^{-1}=\begin{pmatrix}3&-2\\-1&1\end{pmatrix}.
\]
Тогда
\[
A=PDP^{-1}.
\]
Отсюда
\[
A^n=PD^nP^{-1}
=\begin{pmatrix}
3\cdot 2^n-2\cdot 3^n & -2\cdot 2^n+2\cdot 3^n\\
3\cdot 2^n-3\cdot 3^n & -2^{n+1}+3^{n+1}
\end{pmatrix},\qquad n\in\mathbb N.
\]

\item
\[
A=\begin{pmatrix}3&1&0\\0&3&1\\0&0&2\end{pmatrix}.
\]
Так как матрица верхнетреугольная, её собственные значения --- диагональные элементы:
\[
\lambda=3 \quad (\text{кратность }2),\qquad \lambda=2.
\]

Для \(\lambda=3\):
\[
(A-3I)x=0 \iff
\begin{pmatrix}0&1&0\\0&0&1\\0&0&-1\end{pmatrix}x=0
\iff x_2=x_3=0,
\]
то есть
\[
E_3=\langle (1,0,0)^T\rangle.
\]

Для \(\lambda=2\):
\[
(A-2I)x=0 \iff
\begin{pmatrix}1&1&0\\0&1&1\\0&0&0\end{pmatrix}x=0
\iff x_1=-x_2,\ x_3=-x_2,
\]
то есть
\[
E_2=\langle (1,-1,1)^T\rangle.
\]

Геометрическая кратность собственного значения \(3\) равна \(1\), а алгебраическая кратность равна \(2\), поэтому матрица \textbf{не диагонализуема}.

\item \textit{(Проскуряков, №1465)}
\[
A=\begin{pmatrix}
2&-1&2\\
5&-3&3\\
-1&0&-2
\end{pmatrix}.
\]
Характеристический многочлен:
\[
\chi_A(\lambda)=\lambda^3+3\lambda^2+3\lambda+1=(\lambda+1)^3.
\]
Единственное собственное значение:
\[
\lambda=-1 \quad (\text{кратность }3).
\]
Найдём собственные векторы:
\[
(A+I)x=0 \iff
\begin{pmatrix}
3&-1&2\\
5&-2&3\\
-1&0&-1
\end{pmatrix}x=0.
\]
Из последней строки \(x_1=-x_3\). Подставляя, получаем \(x_2=-x_3\). Поэтому
\[
E_{-1}=\langle (-1,-1,1)^T\rangle.
\]

\item \textit{(Проскуряков, №1466)}
\[
A=\begin{pmatrix}
0&1&0\\
-4&4&0\\
-2&1&2
\end{pmatrix}.
\]
Характеристический многочлен:
\[
\chi_A(\lambda)=(\lambda-2)^3.
\]
Единственное собственное значение:
\[
\lambda=2 \quad (\text{кратность }3).
\]
Решим \((A-2I)x=0\):
\[
\begin{pmatrix}
-2&1&0\\
-4&2&0\\
-2&1&0
\end{pmatrix}x=0
\iff x_2=2x_1,
\]
а \(x_3\) --- свободный параметр. Значит,
\[
E_2=\langle (1,2,0)^T,(0,0,1)^T\rangle.
\]

\item \textit{(Проскуряков, №1468)}
\[
A=\begin{pmatrix}
1&-3&3\\
-2&-6&13\\
-1&-4&8
\end{pmatrix}.
\]
Характеристический многочлен:
\[
\chi_A(\lambda)=(\lambda-1)^3.
\]
Следовательно, единственное собственное значение:
\[
\lambda=1 \quad (\text{кратность }3).
\]
Решаем \((A-I)x=0\):
\[
\begin{pmatrix}
0&-3&3\\
-2&-7&13\\
-1&-4&7
\end{pmatrix}x=0
\iff x_1=3x_3,\ x_2=x_3.
\]
Значит,
\[
E_1=\langle (3,1,1)^T\rangle.
\]

\item \textit{(Проскуряков, №1470)}
\[
A=\begin{pmatrix}
7&-12&6\\
10&-19&10\\
12&-24&13
\end{pmatrix}.
\]
Характеристический многочлен:
\[
\chi_A(\lambda)=\lambda^3-\lambda^2-\lambda+1=(\lambda-1)^2(\lambda+1).
\]
Собственные значения:
\[
\lambda=1 \quad (\text{кратность }2),\qquad \lambda=-1.
\]

Для \(\lambda=1\):
\[
(A-I)x=0 \iff
\begin{pmatrix}
6&-12&6\\
10&-20&10\\
12&-24&12
\end{pmatrix}x=0
\iff x_1=2x_2-x_3.
\]
Удобный базис собственного подпространства:
\[
E_1=\langle (2,1,0)^T,\ (-1,0,1)^T\rangle.
\]

Для \(\lambda=-1\):
\[
(A+I)x=0 \iff
\begin{pmatrix}
8&-12&6\\
10&-18&10\\
12&-24&14
\end{pmatrix}x=0.
\]
Получаем, например,
\[
E_{-1}=\langle (3,5,6)^T\rangle.
\]

\item \textit{(Демо 2026, №10)}
Пусть матрица оператора в стандартном базисе равна \(M\). Из условий
\[
M a_1=b_1,\qquad M a_2=b_2,
\]
где
\[
a_1=(2,5)^T,\ a_2=(1,3)^T,\ b_1=(7,-4)^T,\ b_2=(2,-1)^T,
\]
получаем
\[
M\begin{pmatrix}2&1\\5&3\end{pmatrix}
=\begin{pmatrix}7&2\\-4&-1\end{pmatrix}.
\]
Следовательно,
\[
M=\begin{pmatrix}7&2\\-4&-1\end{pmatrix}
\begin{pmatrix}2&1\\5&3\end{pmatrix}^{-1}
=
\begin{pmatrix}11&-3\\-7&2\end{pmatrix}.
\]

\item \textit{(Проскуряков, №1445)}
Докажем существование и единственность линейного преобразования и найдём его матрицу.

Составим матрицы из столбцов-векторов:
\[
A=\begin{pmatrix}
2&0&1\\
3&1&0\\
5&2&0
\end{pmatrix},\qquad
B=\begin{pmatrix}
1&1&2\\
1&1&1\\
1&-1&2
\end{pmatrix}.
\]
Так как
\[
\det A=1\ne 0,
\]
векторы \(a_1,a_2,a_3\) образуют базис пространства, следовательно, существует и единственно линейное преобразование \(\varphi\), переводящее \(a_i\) в \(b_i\).

Его матрица в данном базисе равна
\[
M=BA^{-1}.
\]
Вычисление даёт
\[
M=
\begin{pmatrix}
2&-11&6\\
1&-7&4\\
2&-1&0
\end{pmatrix}.
\]

\item \textit{(Проскуряков, №1446)}
Рассмотрим матрицы, составленные из столбцов-векторов:
\[
A=\begin{pmatrix}
2&4&3\\
0&1&1\\
3&5&2
\end{pmatrix},\qquad
B=\begin{pmatrix}
1&4&1\\
2&5&-1\\
-1&-2&1
\end{pmatrix}.
\]
Так как
\[
\det A=-3\ne 0,
\]
то \(a_1,a_2,a_3\) линейно независимы. Значит, существует и единственно линейное преобразование \(\varphi\), такое что \(\varphi(a_i)=b_i\).

Матрица этого преобразования:
\[
M=BA^{-1}
=
\begin{pmatrix}
-2&\dfrac{11}{3}&\dfrac{5}{3}\\
-4&\dfrac{13}{3}&\dfrac{10}{3}\\
2&-\dfrac{5}{3}&-\dfrac{5}{3}
\end{pmatrix}.
\]

\item \textit{(Демо 2026, №11)}
Пусть исходный базис \(e_1,e_2\), а новый базис \(\hat e_1,\hat e_2\).
Найдём координаты нового базиса в старом.

Из системы
\[
\hat e_1=\begin{pmatrix}4\\3\end{pmatrix},\qquad
\hat e_2=\begin{pmatrix}1\\1\end{pmatrix}
\]
и
\[
e_1=\begin{pmatrix}3\\1\end{pmatrix},\qquad
e_2=\begin{pmatrix}2\\-1\end{pmatrix}
\]
получаем:
\[
\hat e_1=2e_1-e_2,\qquad
\hat e_2=\frac35\,e_1-\frac25\,e_2.
\]
Следовательно, матрица перехода от координат в базисе \((\hat e_1,\hat e_2)\) к координатам в базисе \((e_1,e_2)\) равна
\[
P=\begin{pmatrix}
2&\dfrac35\\
-1&-\dfrac25
\end{pmatrix}.
\]
Если в базисе \(e_1,e_2\) оператор имеет матрицу
\[
A=\begin{pmatrix}
-1&1\\
-3&4
\end{pmatrix},
\]
то в базисе \(\hat e_1,\hat e_2\) его матрица равна
\[
\hat A=P^{-1}AP
=
\begin{pmatrix}
-36&-\dfrac{61}{5}\\
115&39
\end{pmatrix}.
\]

\item
Пусть
\[
A=\begin{pmatrix}1&3\\2&1\end{pmatrix},\qquad
B=\begin{pmatrix}4&5\\0&2\end{pmatrix}.
\]
Тогда
\[
M A=B,\qquad M=BA^{-1}.
\]
Вычисляем:
\[
A^{-1}=\frac1{-5}\begin{pmatrix}1&-3\\-2&1\end{pmatrix},
\qquad
M=\begin{pmatrix}\frac65&\frac75\\[2mm]\frac45&-\frac25\end{pmatrix}.
\]
Это и есть матрица линейного оператора в стандартном базисе.

\item
Пусть
\[
A=\begin{pmatrix}
2&0&1\\
3&1&0\\
5&2&0
\end{pmatrix},\qquad
B=\begin{pmatrix}
1&1&2\\
1&1&1\\
1&-1&2
\end{pmatrix}.
\]
Определитель матрицы \(A\) равен
\[
\det A=1\ne 0,
\]
значит, \(a_1,a_2,a_3\) образуют базис, и искомое преобразование существует и единственно.
Матрица преобразования:
\[
M=BA^{-1}
=\begin{pmatrix}
2&-11&6\\
1&-7&4\\
2&-1&0
\end{pmatrix}.
\]

\item \textit{(Проскуряков, №1452)}
Дана матрица оператора в базисе \(e_1,e_2,e_3,e_4\):
\[
A=\begin{pmatrix}
1&2&0&1\\
3&0&-1&2\\
2&5&3&1\\
1&2&1&3
\end{pmatrix}.
\]

\begin{enumerate}[label=\alph*)]
\item Для базиса \((e_1,e_3,e_2,e_4)\) матрица получается перестановкой 2-го и 3-го базисных векторов:
\[
A_a=\begin{pmatrix}
1&0&2&1\\
2&3&5&1\\
3&-1&0&2\\
1&1&2&3
\end{pmatrix}.
\]

\item Для базиса
\[
f_1=e_1,\quad f_2=e_1+e_2,\quad f_3=e_1+e_2+e_3,\quad f_4=e_1+e_2+e_3+e_4
\]
матрица перехода
\[
P=\begin{pmatrix}
1&1&1&1\\
0&1&1&1\\
0&0&1&1\\
0&0&0&1
\end{pmatrix},
\]
поэтому
\[
A_b=P^{-1}AP=
\begin{pmatrix}
-2&0&1&0\\
1&-4&-8&-7\\
1&4&6&4\\
1&3&4&7
\end{pmatrix}.
\]
\end{enumerate}

\item \textit{(Проскуряков, №1453)}
\[
A=\begin{pmatrix}
15&-11&5\\
20&-15&8\\
8&-7&6
\end{pmatrix},\qquad
f_1=2e_1+3e_2+e_3,\ 
f_2=3e_1+4e_2+e_3,\ 
f_3=e_1+2e_2+2e_3.
\]
Матрица перехода к базису \(f_1,f_2,f_3\):
\[
P=\begin{pmatrix}
2&3&1\\
3&4&2\\
1&1&2
\end{pmatrix},
\qquad
A_f=P^{-1}AP.
\]
Вычисление даёт
\[
A_f=
\begin{pmatrix}
1&0&10\\
0&2&-3\\
0&0&3
\end{pmatrix}.
\]

\item \textit{(Демо 2026, №12)}
\[
A=\begin{pmatrix}
2&-1&1\\
-1&2&-1\\
1&-1&2
\end{pmatrix}.
\]
Матрица симметрична, значит, диагонализуема над \(\mathbb R\), а собственные векторы для разных собственных значений ортогональны.
Характеристический многочлен:
\[
\chi_A(\lambda)=(\lambda-4)(\lambda-1)^2.
\]
Собственные значения:
\[
\lambda=4,\qquad \lambda=1 \ (\text{кратность }2).
\]

Для \(\lambda=4\):
\[
E_4=\langle (1,-1,1)^T\rangle.
\]
Для \(\lambda=1\):
\[
E_1=\langle (1,1,0)^T,\ (-1,0,1)^T\rangle.
\]
Следовательно, базис из собственных векторов существует, например
\[
\bigl((1,-1,1)^T,(1,1,0)^T,(-1,0,1)^T\bigr).
\]

\item
\[
A=\begin{pmatrix}
4&1&0\\
0&4&0\\
0&0&2
\end{pmatrix}.
\]
Собственные значения: \(4\) (кратность \(2\)) и \(2\).

Для \(\lambda=4\):
\[
(A-4I)x=0 \iff x_2=0,\ x_3=0,
\]
то есть
\[
E_4=\langle (1,0,0)^T\rangle.
\]
Для \(\lambda=2\):
\[
E_2=\langle (0,0,1)^T\rangle.
\]
Итак, собственных векторов только два линейно независимых, поэтому базис из собственных векторов построить нельзя.

\item \textit{(Демо 2026, №13)}
Линейный оператор переводит
\[
a_1=(-1,0)^T\mapsto b_1=(5,5)^T,\qquad
a_2=(1,1)^T\mapsto b_2=(-2,-3)^T.
\]
Пусть его матрица \(M\). Тогда
\[
M\begin{pmatrix}-1&1\\0&1\end{pmatrix}
=\begin{pmatrix}5&-2\\5&-3\end{pmatrix},
\]
откуда
\[
M=\begin{pmatrix}-5&3\\-5&2\end{pmatrix}.
\]

\begin{enumerate}[label=\arabic*)]
\item Прямая
\[
2x_1-x_2=-2
\]
параметризуется как \(x=(t,2t+2)\). Тогда
\[
Mx=(t+6,4-t),
\]
и образом будет прямая
\[
x_1+x_2=10.
\]

\item Прообраз этой прямой --- множество решений системы
\[
2(Mx)_1-(Mx)_2=-2.
\]
То есть
\[
2(-5x_1+3x_2)-(-5x_1+2x_2)=-2
\iff -5x_1+4x_2=-2,
\]
или
\[
5x_1-4x_2=2.
\]

\item Прямые, переходящие сами в себя, --- это одномерные инвариантные подпространства, то есть собственные подпространства. У матрицы \(M\) собственные значения комплексные:
\[
\lambda=-\frac32\pm \frac{i\sqrt{11}}2,
\]
поэтому над \(\mathbb R\) инвариантных прямых нет.
\end{enumerate}

\item
Линейный оператор задан:
\[
a_1=(1,0)^T \mapsto b_1=(2,1)^T,\qquad
a_2=(0,1)^T \mapsto b_2=(1,2)^T.
\]
Матрица оператора:
\[
M=\begin{pmatrix}2&1\\1&2\end{pmatrix}.
\]
Образ прямой \(x_1+x_2=1\): положим \(x=(t,1-t)\). Тогда
\[
Mx=(t+1,2-t),
\]
и, исключая \(t\), получаем
\[
y_1+y_2=3.
\]
Значит, образом прямой является прямая \(x_1+x_2=3\).

\item
Найти прообраз прямой
\[
x_1-x_2=1
\]
при линейном операторе
\[
A=\begin{pmatrix}
2&1\\
1&2
\end{pmatrix}.
\]
Имеем
\[
(Ax)_1-(Ax)_2=(2x_1+x_2)-(x_1+2x_2)=x_1-x_2.
\]
Следовательно, прообраз совпадает с самой прямой:
\[
A^{-1}(\{x_1-x_2=1\})=\{x_1-x_2=1\}.
\]

\item
Найти все прямые вида \(\langle v\rangle\), проходящие через начало координат, которые являются инвариантными относительно оператора
\[
A=\begin{pmatrix}
3&1\\
0&2
\end{pmatrix}.
\]
Инвариантные прямые --- это собственные подпространства. Собственные значения: \(\lambda_1=3\), \(\lambda_2=2\).

Для \(\lambda=3\):
\[
E_3=\langle (1,0)^T\rangle.
\]
Для \(\lambda=2\):
\[
(A-2I)x=0 \iff x_1=-x_2,
\]
то есть
\[
E_2=\langle (1,-1)^T\rangle.
\]
Следовательно, искомые прямые:
\[
\langle (1,0)^T\rangle,\qquad \langle (1,-1)^T\rangle.
\]

\item
Линейный оператор переводит
\[
a_1=(1,0)^T\mapsto b_1=(2,1)^T,\qquad
a_2=(0,1)^T\mapsto b_2=(1,2)^T.
\]
Его матрица в стандартном базисе:
\[
M=\begin{pmatrix}2&1\\1&2\end{pmatrix}.
\]
Прямая \(x_1-x_2=1\) переходит в себя, так как
\[
(Mx)_1-(Mx)_2=x_1-x_2.
\]
Значит, образ этой прямой снова есть
\[
x_1-x_2=1.
\]

\item
Линейный оператор переводит векторы
\[
a_1=(1,1)^T,\qquad a_2=(1,-1)^T
\]
соответственно в
\[
b_1=(3,0)^T,\qquad b_2=(0,2)^T.
\]
Тогда
\[
M\begin{pmatrix}1&1\\1&-1\end{pmatrix}
=\begin{pmatrix}3&0\\0&2\end{pmatrix},
\]
откуда
\[
M=\begin{pmatrix}
\frac32&\frac32\\[2mm]
1&-1
\end{pmatrix}.
\]
Собственные значения:
\[
\lambda_1=2,\qquad \lambda_2=-\frac32.
\]
Следовательно, все инвариантные прямые:
\[
\langle (3,1)^T\rangle,\qquad \langle (-1,2)^T\rangle.
\]

\item \textit{(Демо 2026, №14)}
\[
A_\phi=\begin{pmatrix}
1&-1&2&4&-2\\
3&9&-14&2&1\\
3&6&-9&1&1
\end{pmatrix}.
\]
После приведения к ступенчатому виду получаем \(\operatorname{rank}A_\phi=3\).
Значит,
\[
\dim\ker\phi=5-3=2.
\]

Из системы \(A_\phi x=0\) получаем базис ядра:
\[
\ker\phi=\langle (-1,5,3,0,0)^T,\ (-1,-1,0,3,6)^T\rangle.
\]

Базис образа можно взять по опорным столбцам матрицы:
\[
\operatorname{Im}\phi=\langle (1,3,3)^T,\ (-1,9,6)^T,\ (4,2,1)^T\rangle.
\]
Так как \(\operatorname{rank}A_\phi=3=\dim\mathbb R^3\), отображение \textbf{сюръективно}.

\item
\[
A_\phi=\begin{pmatrix}
1&2&-1&0\\
0&1&1&1\\
2&5&0&1
\end{pmatrix}.
\]
Ступенчатый вид даёт ранг \(3\), поэтому
\[
\dim\ker\phi=4-3=1.
\]
Решая \(A_\phi x=0\), получаем
\[
\ker\phi=\langle (2,-1,0,1)^T\rangle.
\]
Опорные столбцы --- первые три, значит,
\[
\operatorname{Im}\phi=\langle (1,0,2)^T,\ (2,1,5)^T,\ (-1,1,0)^T\rangle.
\]
Так как ранг равен \(3\), отображение \textbf{сюръективно}.

\item
\[
A_\psi=\begin{pmatrix}
1&0&2&-1&3\\
0&1&1&0&2\\
1&1&3&-1&5\\
2&1&5&-2&8
\end{pmatrix}.
\]
Ранг матрицы равен \(2\). Значит,
\[
\dim\ker\psi=5-2=3.
\]
Из системы \(A_\psi x=0\) получаем
\[
\ker\psi=\langle (-2,-1,1,0,0)^T,\ (1,0,0,1,0)^T,\ (-3,-2,0,0,1)^T\rangle.
\]
Базис образа можно взять по первым двум столбцам:
\[
\operatorname{Im}\psi=\langle (1,0,1,2)^T,\ (0,1,1,1)^T\rangle.
\]
Так как \(\dim\ker\psi>0\), отображение \textbf{не инъективно}. Поскольку \(\operatorname{rank}\psi=2<4\), оно также \textbf{не сюръективно}.

\item \textit{(Демо 2026, №17)}
Даны векторы
\[
e_1=(0,1,1)^T,\quad e_2=(-1,-1,1)^T,\quad e_3=(1,0,1)^T
\]
в ортонормированном базисе, а матрица оператора \(f\) в базисе \(e_1,e_2,e_3\):
\[
A_f=\begin{pmatrix}
3&1&1\\
1&5&0\\
-3&2&7
\end{pmatrix}.
\]
Пусть
\[
G=\bigl(\langle e_i,e_j\rangle\bigr)=
\begin{pmatrix}
2&0&1\\
0&3&0\\
1&0&2
\end{pmatrix}.
\]
Для матрицы сопряжённого оператора в том же базисе справедливо
\[
A_{f^*}=G^{-1}A_f^T G.
\]
Вычисление даёт
\[
A_{f^*}=
\begin{pmatrix}
-1&2&-7\\
\frac43&5&\frac53\\
5&-1&11
\end{pmatrix}.
\]

\item \textit{(Проскуряков, №1541)}
Пусть \(f_1=e_1,\ f_2=e_1+e_2\). Матрица \(\varphi\) в этом базисе:
\[
A=\begin{pmatrix}1&2\\1&-1\end{pmatrix}.
\]
Так как базис не ортонормирован, используем Gram-матрицу
\[
G=\begin{pmatrix}1&1\\1&2\end{pmatrix}.
\]
Тогда
\[
A_{\varphi^*}=G^{-1}A^TG=
\begin{pmatrix}3&6\\-1&-3\end{pmatrix}.
\]

\item \textit{(Проскуряков, №1542)}
Векторы базиса заданы в ортонормированном базисе:
\[
f_1=(1,2,1)^T,\quad f_2=(1,1,2)^T,\quad f_3=(1,1,0)^T.
\]
Матрица оператора в базисе \(f_1,f_2,f_3\):
\[
A=\begin{pmatrix}
1&1&3\\
0&5&-1\\
2&7&-3
\end{pmatrix}.
\]
Gram-матрица:
\[
G=\begin{pmatrix}
6&5&3\\
5&6&2\\
3&2&2
\end{pmatrix}.
\]
Следовательно,
\[
A_{\varphi^*}=G^{-1}A^TG=
\begin{pmatrix}
-36&-37&-15\\
30&30&14\\
26&27&9
\end{pmatrix}.
\]

\item
Если в некотором ортонормированном базисе
\[
A_f=\begin{pmatrix}
2&1&0\\
1&3&1\\
0&1&2
\end{pmatrix},
\]
то матрица сопряжённого оператора в том же базисе равна транспонированной. Поскольку матрица симметрична, получаем
\[
A_{f^*}=A_f=
\begin{pmatrix}
2&1&0\\
1&3&1\\
0&1&2
\end{pmatrix}.
\]

\item
\[
A=\begin{pmatrix}1&2\\3&4\end{pmatrix}.
\]
Характеристический многочлен:
\[
\chi_A(\lambda)=\lambda^2-5\lambda-2.
\]
Собственные значения:
\[
\lambda_{1,2}=\frac{5\pm\sqrt{33}}2.
\]
Для \(\lambda_1=\frac{5+\sqrt{33}}2\) можно взять собственный вектор
\[
v_1=(4,\,3+\sqrt{33})^T.
\]
Для \(\lambda_2=\frac{5-\sqrt{33}}2\) ---
\[
v_2=(4,\,3-\sqrt{33})^T.
\]
Следовательно,
\[
P=\begin{pmatrix}4&4\\3+\sqrt{33}&3-\sqrt{33}\end{pmatrix},
\qquad
D=\begin{pmatrix}\lambda_1&0\\0&\lambda_2\end{pmatrix},
\qquad
A=PDP^{-1}.
\]

\begin{enumerate}[label=\arabic*)]
\item В базисе собственных векторов матрица оператора равна \(D\).

\item
\[
A^n=PD^nP^{-1},\qquad n\in\mathbb N.
\]
Если раскрыть это произведение, то
\[
A^n=
\begin{pmatrix}
\dfrac{(\sqrt{33}+11)\lambda_2^n+(11-\sqrt{33})\lambda_1^n}{22}
&
\dfrac{2}{\sqrt{33}}\bigl(\lambda_1^n-\lambda_2^n\bigr)
\\[3mm]
\dfrac{3}{\sqrt{33}}\bigl(\lambda_1^n-\lambda_2^n\bigr)
&
\dfrac{(\sqrt{33}-3)\lambda_2^n+(3+\sqrt{33})\lambda_1^n}{2\sqrt{33}}
\end{pmatrix}.
\]

\item Прямая \(x_1=1\) параметризуется как \((1,t)\). Тогда
\[
A\begin{pmatrix}1\\t\end{pmatrix}
=
\begin{pmatrix}1+2t\\3+4t\end{pmatrix}.
\]
Исключая \(t\), получаем уравнение образа:
\[
y_2=2y_1+1.
\]
\end{enumerate}
\end{enumerate}
